\SourceLeaf{192}
adeoque correctio valoris approximati integralis
\[
=-\frac{11}{52500}L^{\mathrm{VI}}-\frac{13}{112500}L^{\mathrm{VIII}}-\frac{5933}{137500000}L^{\mathrm{X}}-\text{etc.}
\]
\SeventhArticle{15}
Coëfficiens $K^{(m)}$ functionis $y$ in seriem evolutae fit, per theorema \textsc{Taylori},
aequalis valori ipsius
\[
\frac{1}{1.2.3\mathinner{.\,.\,.\,.}m}\mathbin{.}\frac{\mathrm{d}^m y}{\mathrm{d}t^m}
\quad\text{sive}\quad
\frac{\Delta^m}{1.2.3\mathinner{.\,.\,.\,.}m}\mathbin{.}\frac{\mathrm{d}^m y}{\mathrm{d}x^m}
\]
pro $t=0$ sive $x=g$; perinde coëfficiens $L^{(m)}$ est valor eiusdem expressionis
pro $t=\frac12$ sive $u=0$ sive $x=g+\frac12\Delta$: utrique coëfficienti \textit{ordinem} $m$ tribuemus. Generaliter itaque loquendo integratio nostra usque ad ordinem $n$ inclus.
exacta erit, quicunque valores pro $a,a',a''\mathinner{.\,.\,.\,.}a^{(n)}$ accipiantur. Attamen hinc
nihil obstat, quominus pro valoribus harum quantitatum scite electis praecisio ad
altiorem gradum evehatur. Ita iam supra vidimus, in methodo \textsc{Cotesii} i. e. pro
$a=0,a'=\frac1n,a''=\frac2n,a'''=\frac3n$ etc. praecisionem sponte ad ordinem $n+1$ inclus.
extendi, quoties $n$ sit numerus par. Generaliter patet, si valores $a,a',a'',a'''$ etc.
ita fuerint electi, ut in functione $T''$ vel $U''$ ab initio excidat terminus unus pluresve, praecisionem totidem gradibus ultra ordinem $n$ promotum iri, quot termini
exciderint. Hinc facile colligitur, quum multitudo quantitatum quas eligere conceditur sit $n+1$, per idoneam earum determinationem praecisionem semper ad
ordinem $2n+1$ inclus. evehi posse, quo pacto adiumento $n+1$ terminorum eundem praecisionis ordinem assequi licebit, ad quem attingendum $2n+1$ vel $2n+2$
terminos in usum vocare oporteret, si \textsc{Cotesii} methodum sequeremur.

\SeventhArticle{16}
Totum hoc negotium in eo vertitur, ut pro quovis valore dato ipsius $n$ functionem $T$ eruamus formae $t^{n+1}+\alpha t^n+\alpha't^{n-1}+\alpha''t^{n-2}$ etc. itaque comparatam, ut
in producto
\[
T(t^{-1}+\tfrac12t^{-2}+\tfrac13t^{-3}+\tfrac14t^{-4}+\text{etc.})
\]
evoluto potestates $t^{-1},t^{-2},t^{-3}\mathinner{.\,.\,.\,.}t^{-(n+1)}$ omnes nanciscantur coëfficientem $0$;
aut si magis placet, functionem $U$ formae $u^{n+1}+\beta u^n+\beta'u^{n-1}+\beta''u^{n-2}+$ etc.,
cuius productum per $u^{-1}+\frac13u^{-3}+\frac15u^{-5}+\frac17u^{-7}+$ etc. liberum evadat a potesta\SourceLeaf{193}%
tibus $u^{-1},u^{-2},u^{-3},u^{-4}\mathinner{.\,.\,.\,.}u^{-(n+1)}$. Modus posterior aliquanto simplicior erit:
quum enim facile perspiciatur, coëfficientes ipsius $U$, ut conditioni praescriptae
satisfiat, alternatim evanescere debere, sive statui $\beta=0,\beta''=0,\beta''''=0$ etc.,
laboris dimidia fere pars iam absoluta censenda erit. Evolvamus casus quosdam
simpliciores.

I. Pro $n=0$, coëfficiens unicus ipsius $t^{-1}$ in producto
\[
(t+\alpha)(t^{-1}+\tfrac12t^{-2}+\tfrac13t^{-3}+\text{etc.})
\]
evanescere debet. Qui quum fiat $=\frac12+\alpha$, habemus $\alpha=-\frac12$, sive $T=t-\frac12$.
Perinde $U=u$.

II. Pro $n=1$, determinatio ipsius $T$ pendet a duabus aequationibus
\[
\begin{aligned}
0&=\tfrac13+\tfrac12\alpha+\alpha'\\
0&=\tfrac14+\tfrac13\alpha+\tfrac12\alpha'
\end{aligned}
\]
unde deducimus $\alpha=-1,\alpha'=+\frac16$, sive $T=tt-t+\frac16$. Determinatio functionis $U$ unicam aequationem affert
\[
0=\tfrac13+\beta'
\]
unde $\beta'=-\frac13$, sive $U=uu-\frac13$.

III. Pro $n=2$, functio $T$ determinatur adiumento trium aequationum
\[
\begin{aligned}
0&=\tfrac14+\tfrac13\alpha+\tfrac12\alpha'+\alpha''\\
0&=\tfrac15+\tfrac14\alpha+\tfrac13\alpha'+\tfrac12\alpha''\\
0&=\tfrac16+\tfrac15\alpha+\tfrac14\alpha'+\tfrac13\alpha''
\end{aligned}
\]
unde nanciscimur $\alpha=-\frac32,\alpha'=\frac35,\alpha''=-\frac1{20}$, adeoque $T=t^3-\frac32tt+\frac35t-\frac1{20}$.
Ad determinandam $U$ unica aequatio sufficit
\[
0=\tfrac15+\tfrac13\beta'
\]
unde $\beta'=-\frac35$ sive $U=u^3-\frac35u$.

Attamen hunc modum, qui calculos continuo molestiores adducit, hic ulterius non persequemur, sed ad fontem genuinum solutionis generalis progrediemur.
\SourceLeaf{194}
\SeventhArticle{17}
Proposita fractione continua
\[
\varphi=\cfrac{v}{w+\cfrac{v'}{w'+\cfrac{v''}{w''+\cfrac{v'''}{w'''+\text{etc.}}}}}
\]
constat, fractiones continuo magis appropinquantes inveniri per algorithmum
sequentem. Formentur duae quantitatum series, $V,V',V'',V'''$ etc.,
$W,W',W'',W'''$ etc. per hasce formulas
\[
\begin{array}{ll}
V=0&W=1\\[4pt]
V'=v&W'=wW\\[4pt]
V''=w'V'+v'V&W''=w'W'+v'W\\[4pt]
V'''=w''V''+v''V'&W'''=w''W''+v''W'\\[4pt]
V''''=w'''V'''+v'''V''&W''''=w'''W'''+v'''W''
\end{array}
\]
etc. eritque
\[
\begin{aligned}
\cfrac VW&=0\\[9pt]
\cfrac{V'}{W'}&=\cfrac vw\\[9pt]
\cfrac{V''}{W''}&=\cfrac v{w+\cfrac{v'}{w'}}\\[9pt]
\cfrac{V'''}{W'''}&=\cfrac v{w+\cfrac{v'}{w'+\cfrac{v''}{w''}}}
\end{aligned}
\]
et sic porro. Praeterea constat, vel facile ex ipsis aequationibus praecedentibus
confirmatur, esse
\[
\begin{aligned}
VW'-V'W&=-v\\
V'W''-V''W'&=+vv'\\
V''W'''-V'''W''&=-vv'v''\\
V'''W''''-V''''W'''&=+vv'v''v'''
\end{aligned}
\]
etc. Hinc perspicuum est, seriei
\[
\frac v{WW'}-\frac{vv'}{W'W''}+\frac{vv'v''}{W''W'''}-\frac{vv'v''v'''}{W'''W''''}+\text{etc.}
\]
\SourceLeaf{195}
\[
\begin{array}{ll}
\text{terminum primum esse}&=\dfrac{V'}{W'}\\[10pt]
\text{summam duorum terminorum primorum}&=\dfrac{V''}{W''}\\[10pt]
\text{summam trium terminorum primorum}&=\dfrac{V'''}{W'''}\\[10pt]
\text{summam quatuor terminorum primorum}&=\dfrac{V''''}{W''''}
\end{array}
\]
et sic porro; quocirca series ipsa vel in infinitum vel usque dum abrumpatur continuata ipsam fractionem continuam $\varphi$ exprimet. Simul hinc habetur differentia
inter $\varphi$ atque singulas fractiones appropinquantes $\frac{V'}{W'},\frac{V''}{W''},\frac{V'''}{W'''}$ etc.

E formula 33 art. 14 \textit{Disquisitionum generalium circa seriem infinitam} mutando $t$ in $\frac1u$, facile obtinemus transformationem seriei
\[
\varphi=u^{-1}+\tfrac13u^{-3}+\tfrac15u^{-5}+\tfrac17u^{-7}+\text{etc.}
\]
in fractionem continuam sequentem
\[
\cfrac{1}{u-\cfrac{\cfrac13}{u-\cfrac{\cfrac{2\mathbin{.}2}{3\mathbin{.}5}}{u-\cfrac{\cfrac{3\mathbin{.}3}{5\mathbin{.}7}}{u-\cfrac{\cfrac{4\mathbin{.}4}{7\mathbin{.}9}}{u-\text{etc.}}}}}}
\]
ita ut habeatur
\[
\begin{aligned}
v&=1,\ v'=-\tfrac13,\ v''=-\tfrac4{15},\ v'''=-\tfrac9{35},\ v''''=-\tfrac{16}{63}\ \text{etc.}\\
w&=w'=w''=w'''=w''''\ \text{etc.}=u.
\end{aligned}
\]
Hinc pro $V,V',V'',V'''$ etc. $W,W',W'',W'''$ etc. nanciscimur valores sequentes
\[
\begin{aligned}&V=0,\\[5pt]&\qquad W=1\\[5pt]&V'=1,\\[5pt]&\qquad W'=u\\[5pt]&V''=u,\\[5pt]&\qquad W''=uu-\tfrac13\\[5pt]&V'''=uu-\tfrac4{15},\\[5pt]&\qquad W'''=u^3-\tfrac35u\end{aligned}\]
\[\begin{aligned}&V''''=u^3-\tfrac{11}{21}u,\\[5pt]&\qquad W''''=u^4-\tfrac67uu+\tfrac3{35}\\[5pt]&V^{\mathrm{V}}=u^4-\tfrac79uu+\tfrac{64}{945},\\[5pt]&\qquad W^{\mathrm{V}}=u^5-\tfrac{10}{9}u^3+\tfrac5{21}u\\[5pt]&V^{\mathrm{VI}}=u^5-\tfrac{34}{33}u^3+\tfrac15u,\\[5pt]&\qquad W^{\mathrm{VI}}=u^6-\tfrac{15}{11}u^4+\tfrac5{11}uu-\tfrac5{231}\\[5pt]&V^{\mathrm{VII}}=u^6-\tfrac{50}{39}u^4+\tfrac{283}{715}uu-\tfrac{256}{15015},\\[5pt]&\qquad W^{\mathrm{VII}}=u^7-\tfrac{21}{13}u^5+\tfrac{105}{143}u^3-\tfrac{35}{429}u\ \text{etc.}\end{aligned}
\]
\SourceLeaf{196}
Levi attentione adhibita elucet, singulas $V,V',V'',V'''$ etc. $W,W',W'',W'''$ etc.
fieri functiones integras indeterminatae $u$; terminum altissimum in $V^{(m)}$ fieri
$u^{m-1}$, potestatesque $u^{m-2},u^{m-4},u^{m-6}$ etc. abesse; terminum altissimum vero in
$W^{(m)}$ fieri $u^m$, atque abesse potestates $u^{m-1},u^{m-3},u^{m-5}$ etc. Per ea autem, quae
supra demonstravimus, erit
\[
\begin{aligned}&\varphi=\frac1{WW'}+\frac1{3W'W''}+\frac{2\mathbin{.}2}{3\mathbin{.}3\mathbin{.}5W''W'''}\\[5pt]&+\frac{2\mathbin{.}2\mathbin{.}3\mathbin{.}3}{3\mathbin{.}3\mathbin{.}5\mathbin{.}5\mathbin{.}7W'''W''''}+\frac{2\mathbin{.}2\mathbin{.}3\mathbin{.}3\mathbin{.}4\mathbin{.}4}{3\mathbin{.}3\mathbin{.}5\mathbin{.}5\mathbin{.}7\mathbin{.}7\mathbin{.}9W''''W'''''}+\text{etc.}\end{aligned}
\]
ac proin generaliter
\[
\begin{aligned}
\varphi-\frac{V^{(m)}}{W^{(m)}}&=\frac{2\mathbin{.}2\mathbin{.}3\mathbin{.}3\mathinner{.\,.\,.\,.\,.} m\mathbin{.}m}{3\mathbin{.}3\mathbin{.}5\mathbin{.}5\mathinner{.\,.\,.\,.}(2m-1)(2m+1)W^{(m)}W^{(m+1)}}\\[5pt]
&\quad+\frac{2\mathbin{.}2\mathbin{.}3\mathbin{.}3\mathinner{.\,.\,.\,.\,.}(m+1)(m+1)}{3\mathbin{.}3\mathbin{.}5\mathbin{.}5\mathinner{.\,.\,.\,.\,.}(2m+1)(2m+3)W^{(m+1)}W^{(m+2)}}\\[5pt]
&\quad+\text{etc.}
\end{aligned}
\]
Si igitur $\varphi-\frac{V^{(m)}}{W^{(m)}}$ in seriem descendentem convertitur, eius terminus primus erit
\[
=\frac{2\mathbin{.}2\mathbin{.}3\mathbin{.}3\mathinner{.\,.\,.\,.} m\mathbin{.}m u^{-(2m+1)}}{3\mathbin{.}3\mathbin{.}5\mathbin{.}5\mathinner{.\,.\,.\,.}(2m-1)(2m+1)}
\]
Productum vero $\varphi W^{(m)}$ compositum erit e functione integra $V^{(m)}$ atque serie
infinita, cuius terminus primus
\[
=\frac{2\mathbin{.}2\mathbin{.}3\mathbin{.}3\mathinner{.\,.\,.\,.\,.} mmu^{-(m+1)}}{3\mathbin{.}3\mathbin{.}5\mathbin{.}5\mathinner{.\,.\,.\,.\,.}(2m-1)(2m+1)}
\]
Hinc igitur sponte inventa est functio $U$ ordinis $n+1$; quae conditioni in
art. praec. stabilitae satisfacit, scilicet ut productum $\varphi U$ liberum evadat a potestatibus $u^{-1},u^{-2},u^{-3}\mathinner{.\,.\,.\,.\,.} u^{-(n+1)}$. Scilicet non est alia quam $W^{(n+1)}$, simulque patet, $U'$ aequalem fieri ipsi $V^{(m+1)}$, nec non terminum primum ipsius
$U''$ esse
\[
=\frac{2\mathbin{.}2\mathbin{.}3\mathbin{.}3\mathinner{.\,.\,.\,.\,.}(n+1)(n+1)}{3\mathbin{.}3\mathbin{.}5\mathbin{.}5\mathinner{.\,.\,.\,.\,.}(2n+1)(2n+3)}\mathbin{.}u^{-(n+2)}
\]
Quodsi igitur pro $b,b',b''\mathinner{.\,.\,.\,.\,.} b^{(n)}$ accipiuntur radices aequationis $W^{(n+1)}=0$,
valoresque coëfficientium $R,R',R''\mathinner{.\,.\,.\,.\,.} R^{(n)}$ per praecepta supra tradita eruntur, formula nostra integralis praecisione gaudebit ad ordinem $2n+1$ ascendente, eiusque correctio exprimetur proxime per
\[
\begin{aligned}&\frac1{2^{2n+2}}\mathbin{.}\frac{2\mathbin{.}2\mathbin{.}3\mathbin{.}3\mathinner{.\,.\,.\,.\,.}(n+1)(n+1)}{3\mathbin{.}3\mathbin{.}5\mathbin{.}5\mathinner{.\,.\,.\,.\,.}(2n+1)(2n+3)}L^{(2n+2)}\\[5pt]&=\frac{1\mathbin{.}1\mathbin{.}2\mathbin{.}2\mathbin{.}3\mathbin{.}3\mathinner{.\,.\,.\,.\,.}(n+1)(n+1)}{2\mathbin{.}6\mathbin{.}6\mathbin{.}10\mathbin{.}10\mathbin{.}14\mathinner{.\,.\,.\,.\,.}(4n+2)(4n+6)}L^{(2n+2)}\end{aligned}
\]
\SourceLeaf{197}
\SeventhArticle{18}
Disquisitiones art. praec. functiones idoneas $U$ pro singulis valoribus numeri $n$ invenire quidem docent, sed successive tantum, dum a valoribus minoribus ad maiores transeundum est. Facile autem animadvertimus, has functiones
generaliter exprimi per
\[
\begin{aligned}&
u^{n+1}-\frac{(n+1)n}{2\mathbin{.}(2n+1)}u^{n-1}
\\[5pt]&+\frac{(n+1)n(n-1)(n-2)}{2\mathbin{.}4(2n+1)(2n-1)}u^{n-3}
\\[5pt]&-\frac{(n+1)n(n-1)(n-2)(n-3)(n-4)}{2\mathbin{.}4\mathbin{.}6\mathbin{.}(2n+1)(2n-1)(2n-3)}u^{n-5}\\ &
+\text{etc.}
\end{aligned}
\]
sive etiam, si characteristica $F$ ad normam commentationis supra citatae utimur, per
\[
u^{n+1}F(-\tfrac12n,\ -\tfrac12(n+1),\ -(n+\tfrac12),\ u^{-2})
\]
Haecce inductio facile in demonstrationem rigorosam convertitur per methodum
vulgo notam, aut, si ita videtur, adiumento formulae 19 in comment. cit. Functio
$U$, si magis placet, etiam ordine terminorum inverso, exprimi potest per
\[
\pm\frac{3\mathbin{.}5\mathbin{.}7\mathinner{.\,.\,.\,.}(n+1)}{(n+3)(n+5)\mathinner{.\,.\,.\,.}(2n+1)}\mathbin{.}uF(-\tfrac12n,\ \tfrac12(n+3),\ \tfrac32,\ uu)
\]
pro $n$ pari, valente signo superiori vel inferiori, prout $\frac12n$ par est vel impar
aut per
\[
\pm\frac{1\mathbin{.}3\mathbin{.}5\mathinner{.\,.\,.\,.} n}{(n+2)(n+4)\mathinner{.\,.\,.\,.}(2n+1)}F(-\tfrac12(n+1),\ \tfrac12n+1,\ \tfrac12,\ uu)
\]
pro $n$ impari, valente signo superiori vel inferiori, prout $\frac12(n+1)$ par est vel impar.

Functio $U'$ expressionem generalem aeque simplicem non admittit: facile
tamen ex ipsa genesi quantitatum $V,V',V''$ etc. colligitur, terminum ultimum
ipsius $U'$ pro $n$ pari fieri
\[
=\pm\frac{2\mathbin{.}2\mathbin{.}4\mathbin{.}4\mathbin{.}6\mathbin{.}6\mathinner{.\,.\,.\,.\,.} n\mathbin{.}n}{3\mathbin{.}5\mathbin{.}7\mathbin{.}9\mathbin{.}11\mathbin{.}13\mathinner{.\,.\,.\,.\,.}(2n-1)(2n+1)}
\]
signo superiori vel inferiori valente, prout $\frac12n$ par est vel impar.

Functio $U''=\varphi W^{(n+1)}-V^{(n+1)}$, cuius terminum primum iam in art.
praec. assignare docuimus, etiam per algorithmum recurrentem evolvi potest,
quum manifesto generaliter habeatur
\[
\begin{aligned}
\varphi W''-V''&=w'(\varphi W'-V')+v'(\varphi W-V)\\
\varphi W'''-V'''&=w''(\varphi W''-V'')+v''(\varphi W'-V')\\
\varphi W''''-V''''&=w'''(\varphi W'''-V''')+v'''(\varphi W''-V'')
\end{aligned}
\]
\SourceLeaf{198}
etc. adeoque eo quem tractamus casu
\[
\begin{aligned}&\varphi W^{(m+2)}-V^{(m+2)}=u(\varphi W^{(m+1)}-V^{(m+1)})\\[5pt]&-\frac{(m+1)^2}{(2m-1)(2m+1)}(\varphi W^{(m)}-V^{(m)})\end{aligned}
\]
Ita invenimus
\[
\begin{aligned}
\varphi W-V&=u^{-1}+\tfrac13u^{-3}+\tfrac15u^{-5}+\tfrac17u^{-7}+\text{etc.}\\
\varphi W'-V'&=\tfrac13u^{-2}+\tfrac15u^{-4}+\tfrac17u^{-6}+\tfrac19u^{-8}+\text{etc.}\\
\varphi W''-V''&=\tfrac4{45}u^{-3}+\tfrac8{105}u^{-5}+\tfrac4{63}u^{-7}+\tfrac{112}{2079}u^{-9}+\text{etc.}\\
\varphi W'''-V'''&=\tfrac4{175}u^{-4}+\tfrac8{315}u^{-6}+\tfrac4{165}u^{-8}+\tfrac{16}{715}u^{-10}+\text{etc.}
\end{aligned}
\]
etc. quas series ita quoque exhibere licet
\[
\begin{aligned}
\varphi W-V&=u^{-1}\left(1+\frac{1\mathbin{.}2}{2\mathbin{.}3}u^{-2}\right.\\[3pt]&\qquad\left.+\frac{1\mathbin{.}2\mathbin{.}3\mathbin{.}4}{2\mathbin{.}4\mathbin{.}3\mathbin{.}5}u^{-4}+\frac{1\mathbin{.}2\mathbin{.}3\mathbin{.}4\mathbin{.}5\mathbin{.}6}{2\mathbin{.}4\mathbin{.}6\mathbin{.}3\mathbin{.}5\mathbin{.}7}u^{-6}+\text{etc.}\right)\\[4pt]
\varphi W'-V'&=\tfrac13u^{-2}\left(1+\frac{2\mathbin{.}3}{2\mathbin{.}5}u^{-4}\right.\\[3pt]&\qquad\left.+\frac{2\mathbin{.}3\mathbin{.}4\mathbin{.}5}{2\mathbin{.}4\mathbin{.}5\mathbin{.}7}u^{-4}+\frac{2\mathbin{.}3\mathbin{.}4\mathbin{.}5\mathbin{.}6\mathbin{.}7}{2\mathbin{.}4\mathbin{.}6\mathbin{.}5\mathbin{.}7\mathbin{.}9}u^{-6}+\text{etc.}\right)\\[4pt]
\varphi W''-V''&=\tfrac4{45}u^{-3}\left(1+\frac{3\mathbin{.}4}{2\mathbin{.}7}u^{-2}\right.\\[3pt]&\qquad\left.+\frac{3\mathbin{.}4\mathbin{.}5\mathbin{.}6}{2\mathbin{.}4\mathbin{.}7\mathbin{.}9}u^{-4}+\frac{3\mathbin{.}4\mathbin{.}5\mathbin{.}6\mathbin{.}7\mathbin{.}8}{2\mathbin{.}4\mathbin{.}6\mathbin{.}7\mathbin{.}9\mathbin{.}11}u^{-6}+\text{etc.}\right)\\[4pt]
\varphi W'''-V'''&=\tfrac4{175}u^{-4}\left(1+\frac{4\mathbin{.}5}{2\mathbin{.}9}u^{-2}\right.\\[3pt]&\qquad\left.+\frac{4\mathbin{.}5\mathbin{.}6\mathbin{.}7}{2\mathbin{.}4\mathbin{.}9\mathbin{.}11}u^{-4}+\frac{4\mathbin{.}5\mathbin{.}6\mathbin{.}7\mathbin{.}8\mathbin{.}9}{2\mathbin{.}4\mathbin{.}6\mathbin{.}9\mathbin{.}11\mathbin{.}13}u^{-6}+\text{etc.}\right)
\end{aligned}
\]
etc. Hanc inductionem sequentes habebimus generaliter
\[
U''=\varphi W^{(n+1)}-V^{(n+1)}\quad\text{aequalem producto ex}
\]
\[
\frac{2\mathbin{.}2\mathbin{.}3\mathbin{.}3\mathbin{.}4\mathbin{.}4\mathinner{.\,.\,.\,.\,.}(n+1)\mathbin{.}(n+1)}{3\mathbin{.}3\mathbin{.}5\mathbin{.}5\mathbin{.}7\mathbin{.}7\mathbin{.}9\mathinner{.\,.\,.\,.\,.}(2n+1)(2n+3)}u^{-(n+2)}
\]
in seriem infinitam
\[
1+\frac{(n+2)(n+3)}{2(2n+5)}u^{-2}+\frac{(n+2)(n+3)(n+4)(n+5)}{2\mathbin{.}4\mathbin{.}(2n+5)(2n+7)}u^{-4}+\text{etc.}
\]
aut si magis placet in $F(\frac12n+1,\ \frac12n+\frac32,\ n+\frac52,\ u^{-2})$. Haec quoque inductio
facillime ad plenam certitudinem evehitur vel per methodum vulgo notam vel adiumento formulae 19 in commentatione saepius citatae.

\SeventhArticle{19}
Quum sufficiat, functionum $T,U$ alterutram nosse, posterioris determinationem tamquam simpliciorem praetulimus. Quae quemadmodum evolutioni
seriei $u^{-1}+\frac13u^{-3}+\frac15u^{-5}+\text{etc.}$ in fractionem continuam innixa est, per ratiocinia similia ex evolutione seriei $t^{-1}+\frac12t^{-2}+\frac13t^{-3}+\frac14t^{-4}+\text{etc.}$ in fractionem
continuam
\SourceLeaf{199}
\[
\cfrac1{t-\cfrac{\frac12}{1-\cfrac{\frac16}
{t-\cfrac{\frac26}{1-\cfrac{\frac2{10}}
{t-\cfrac{\frac3{10}}{1-\mathrm{etc.}}}}}}}
\]
derivare potuissemus algorithmum ad determinandam
functionem $T$ pro valoribus successivis numeri $n$. Ad eandem
vero conclusionem pervenimus perpendendo, $T$ nihil aliud esse
quam $\frac U{2^{n+1}}$ seu $\frac{W^{(n+1)}}{2^{n+1}}$,
si pro $u$ scribitur $2t-1$, quo pacto functiones successive pro
$T$ adoptandae habentur per algorithmum sequentem:
\[
\begin{aligned}
W&=1\\
\tfrac12W'&=t-\tfrac12\\
\tfrac14W''&=(t-\tfrac12).\tfrac12W'-\frac{1.1}{2.6}W
\\[3pt]&\qquad=tt-t+\tfrac16\\
\tfrac18W'''&=(t-\tfrac12).\tfrac14W''-\frac{2.2}{6.10}.\tfrac12W'
\\[3pt]&\qquad=t^3-\tfrac32tt+\tfrac35t-\tfrac1{20}\\
\tfrac1{10}W''''&=(t-\tfrac12).\tfrac18W'''-\frac{3.3}{10.14}.\tfrac14W''
\\[3pt]&\qquad=t^4-2t^3+\tfrac97tt-\tfrac27t+\tfrac1{70}
\end{aligned}
\]
etc. Per inductionem hinc resultat generaliter
\[
\begin{aligned}
T={}&t^{n+1}-\frac{(n+1)^2}{1.(2n+2)}t^n
+\frac{(n+1)^2.nn}{1.2.(2n+2)(2n+1)}t^{n-1}\\
&-\frac{(n+1)^2.nn.(n-1)^2}{1.2.3.(2n+2)(2n+1).2n}t^{n-2}
+\mathrm{etc.}
\end{aligned}
\]
sive $T=t^{n+1}F(-(n+1),-(n+1),-2(n+1),t^{-1})$,
cui inductioni facile est demonstrationis vim conciliare.
Si magis arridet, $T$ ordine terminorum inverso etiam per
\[
\pm\frac{1.2.3.4\mathinner{.\,.\,.\,.\,.}(n+1)}
{2.6.10.14\mathinner{.\,.\,.\,.\,.}(4n+2)}
F(n+2,-(n+1),1,t)
\]
exprimi potest, ubi signum superius valet pro $n$ impari,
inferius pro pari. Simili denique modo generaliter $T''$
aequalis invenitur producto ex
\[
\frac{1.1.2.2.3.3\mathinner{.\,.\,.\,.\,.}(n+1).(n+1)}
{2.6.6.10.10.14\mathinner{.\,.\,.\,.\,.}(4n+2).(4n+6)}t^{-(n+2)}
\]
in seriem infinitam
\[
\begin{aligned}
1&+\frac{(n+2)^2}{1.(2n+4)}t^{-1}
+\frac{(n+2)^2(n+3)^2}{1.2.(2n+4)(2n+5)}t^{-2}\\
&+\frac{(n+2)^2.(n+3)^2(n+4)^2}
{1.2.3.(2n+4)(2n+5)(2n+6)}t^{-3}+\mathrm{etc.}
\end{aligned}
\]
sive in $F(n+2,n+2,2n+4,t^{-1})$
\SourceLeaf{200}
\SeventhArticle{20}
Quum in functione $U$ potestates $u^n$, $u^{n-2}$, $u^{n-4}$
etc. absint, e radicibus aequationis $U=0$ binae semper erunt
magnitudine aequales signis oppositae, quibus pro valore pari
ipsius $n$ adhuc associare oportet radicem singularem $0$.
Inventis radicibus, valores coëfficientium $R,R',R''$ etc.
secundum methodum art. 11 habebuntur per functionem integram
ipsius $u$, quae pro valore impari ipsius $n$ erit formae
\[
\gamma u^{n-1}+\gamma'u^{n-3}+\gamma''u^{n-5}+\mathrm{etc.}
\]
pro valore pari autem, si excluditur coëfficiens radici $u=0$
respondens, formae
\[
\gamma u^{n-2}+\gamma'u^{n-4}+\gamma''u^{n-6}+\mathrm{etc.}
\]
Exemplum art. 12 ipsam hanc reductionem exhibet pro $n=6$.
Manifesto igitur valoribus oppositis ipsius $u$ semper respondent
coëfficientes aequales. Ceterum in casu eo, ubi $n$ est par,
coëfficiens radici $u=0$ respondens facile generaliter a priori
assignari potest. Habebitur hic coëfficiens, si in
$\frac{U'}{(\frac{\mathrm{d}U}{\mathrm{d}u})}$ substituitur
$u=0$. Valorem numeratoris $U'$ pro $u=0$ iam in art. 18
tradidimus, valor denominatoris autem ibinde erit
\[
\begin{aligned}&=\pm\frac{3.5.7\mathinner{.\,.\,.\,.\,.}(n+1)}
{(n+3)(n+5)\mathinner{.\,.\,.\,.\,.}(2n+1)}\\[5pt]&=\pm\frac{3.3.5.5.7.7\mathinner{.\,.\,.\,.\,.}(n+1)(n+1)}
{3.5.7.9.11\mathinner{.\,.\,.\,.\,.}(2n+1)}\end{aligned}
\]
adeoque coëfficiens quaesitus
\[
=\left(\frac{2.4.6.8\mathinner{.\,.\,.\,.\,.} n}
{3.5.7.9\mathinner{.\,.\,.\,.\,.}(n+1)}\right)^2
\]

\SeventhArticle{21}
Functio integra ipsius $u$ coëfficientes $R,R',R''$ etc.
repraesentans in eo quem hic tractamus casu etiam independenter
a methodo generali art. 11 erui potest sequenti modo.
Differentiando aequationem
\[
\varphi-\frac{U'}U=\frac{U''}U
\]
substituendo dein $\frac{\mathrm{d}\varphi}{\mathrm{d}u}
=\frac1{1-uu}$, ac multiplicando per $UU(uu-1)$, obtinemus
\[
\begin{aligned}&(uu-1)U'\frac{\mathrm{d}U}{\mathrm{d}u}
-U\left(\frac{\mathrm{d}U'}{\mathrm{d}u}.(uu-1)+U\right)\\[5pt]&=(uu-1)UU\frac{\mathrm{d}(\frac{U''}U)}{\mathrm{d}u}\end{aligned}
\]
\SourceLeaf{201}
Termini huius aequationis ad laevam manifesto constituunt
functionem integram ipsius $u$: itaque necessario in parte ad
dextram coëfficientes potestatum ipsius $u$ cum exponentibus
negativis sese destruere debent.

Sed $\frac{\mathrm{d}\frac{U''}U}{\mathrm{d}u}$ producit
seriem infinitam incipientem a termino
\[
-\left(\frac{1.2.3.4\mathinner{.\,.\,.\,.\,.}(n+1)}
{1.3.5.7\mathinner{.\,.\,.\,.\,.}(2n+1)}\right)^2u^{-(2n+4)}
\]
qua igitur per $(uu-1)UU$ multiplicata nihil aliud prodire
poterit nisi quantitas constans
\[
-\left(\frac{1.2.3.4\mathinner{.\,.\,.\,.\,.}(n+1)}
{1.3.5.7\mathinner{.\,.\,.\,.\,.}(2n+1)}\right)^2
\]
Hinc colligimus\HistoricalNote{20}
\[
(uu-1)U'\frac{\mathrm{d}U}{\mathrm{d}u}
+\left(\frac{1.2.3.4\mathinner{.\,.\,.\,.\,.}(n+1)}
{1.3.5.7\mathinner{.\,.\,.\,.\,.}(2n+1)}\right)^2
\]
divisibilem esse per $U$, quamobrem functioni fractae
$\frac{U'}{(\frac{\mathrm{d}U}{\mathrm{d}u})}$,
quae coëfficientes $R,R',R''$ etc. suggerit,
aequivalebit functio integra
\[
-\left(\frac{1.3.5.7\mathinner{.\,.\,.\,.\,.}(2n+1)}
{1.2.3.4\mathinner{.\,.\,.\,.\,.}(n+1)}U'\right)^2.(uu-1)
\]
Loco huius functionis, quae est ordinis $2n+2$, manifestoque
solas potestates pares ipsius $u$ implicat, adoptari poterit
residuum ex eius divisione per $U$ ortum, quod erit ordinis $n$,
seu $n-1$, prout $n$ par est seu impar. Si vero in casu priori
coëfficientem eum, qui respondet radici $u=0$, excludere malumus,
loco illius functionis eius residuum ex divisione per $\frac Uu$
ortum adoptabimus, quod tantummodo ad ordinem $n-2$ ascendet.

\SeventhArticle{22}
Ut praesto sint, quae ad applicationem methodi hucusque expositae
requiruntur. adiungere visum est, pro valoribus successivis
numeri $n$, valores numericos tum quantitatum $a,a',a''$ etc.,
tum coëfficientium $R,R'.R''$ etc. ad sedecim figuras
computatos, una cum horum logarithmis ad decem figuras.
